\documentclass[12pt, a4paper]{article}

\usepackage[margin=1in]{geometry}
\usepackage{amsmath, amssymb, amsfonts}
\usepackage{physics}
\usepackage{titlesec}
\usepackage{fancyhdr}
\usepackage{enumitem}
\usepackage{hyperref}
\usepackage{xcolor}
\usepackage{tcolorbox}
\tcbuselibrary{skins,breakable}

\hypersetup{
    colorlinks=true,
    linkcolor=blue,
    urlcolor=blue,
    citecolor=blue
}

\setlength{\headheight}{14.5pt}
\pagestyle{fancy}
\fancyhf{}
\fancyhead[L]{24BSE2113-D: PDE, Transforms \& Optimization}
\fancyhead[R]{Phase 1: Theoretical Assignment}
\fancyfoot[C]{\thepage}
\renewcommand{\headrulewidth}{0.4pt}

\titleformat{\section}{\normalfont\Large\bfseries}{\thesection}{1em}{}
\titleformat{\subsection}{\normalfont\large\bfseries}{\thesubsection}{1em}{}

\newtcolorbox{resultbox}{
    colback=blue!5!white,
    colframe=blue!40!black,
    boxrule=0.6pt,
    arc=2pt,
    left=6pt, right=6pt, top=4pt, bottom=4pt
}

\newcommand{\groupNumber}{13}
\newcommand{\studentOneName}{Chriss Mathew}
\newcommand{\studentOneRoll}{23}
\newcommand{\studentTwoName}{Rajan Devika S}
\newcommand{\studentTwoRoll}{25}
\newcommand{\studentThreeName}{Liju Renjitha Babu}
\newcommand{\studentThreeRoll}{55}
\newcommand{\studentFourName}{Sethuparvathy K J}
\newcommand{\studentFourRoll}{58}
\newcommand{\studentFiveName}{Sharon Thomas}
\newcommand{\studentFiveRoll}{61}

\begin{document}

\begin{titlepage}
    \centering
    \vspace*{1cm}

    {\large \textbf{Saintgits College of Engineering (Autonomous)}\par}
    {\large Department of Electronics and Communication Engineering\par}
    \vspace{0.3cm}
    \rule{\linewidth}{0.5pt}
    \vspace{0.5cm}

    {\Large \textbf{Course: 24BSE2113-D}\par}
    {\large Partial Differential Equations, Transforms \& Optimization Techniques\par}
    \vspace{1cm}

    {\Large \textbf{Group Assignment \& Micro-Project-1}\par}
    \vspace{0.2cm}
    {\Large \textbf{Phase 1: Theoretical Assignment}\par}
    \vspace{1cm}

    {\Large \textit{Topic 1: Quantum Confinement}\par}
    {\Large \textit{(The Schr\"odinger Wave Equation)}\par}

    \vfill

    \begin{center}
    \begin{tabular}{ll}
        \textbf{Group Number:} & \groupNumber \\[4pt]
        \textbf{Student 1:} & \studentOneName \; (Roll No: \studentOneRoll) \\[2pt]
        \textbf{Student 2:} & \studentTwoName \; (Roll No: \studentTwoRoll) \\[2pt]
        \textbf{Student 3:} & \studentThreeName \; (Roll No: \studentThreeRoll) \\[2pt]
        \textbf{Student 4:} & \studentFourName \; (Roll No: \studentFourRoll) \\[2pt]
        \textbf{Student 5:} & \studentFiveName \; (Roll No: \studentFiveRoll) \\
    \end{tabular}
    \end{center}

    \vfill
    {\large \today\par}
\end{titlepage}

\section{Introduction}
Quantum confinement is one of the defining phenomena of modern nanoscale
electronics. When the physical dimensions of a semiconductor structure
shrink down to the scale of a few nanometers -- comparable to the
electron's own de Broglie wavelength -- the electron can no longer be
treated as a free classical particle. Instead, it must be described as
a matter wave whose behaviour is governed by quantum mechanics. This
confinement forces the electron's allowed energies to become discrete
rather than continuous, which is the physical basis for technologies
such as quantum dots, quantum well lasers, and modern transistor
channels.

The simplest and most instructive model of this effect is the
\emph{particle in an infinite one-dimensional potential well}: an
electron trapped between two infinitely high potential barriers a
distance $L$ apart. Despite its simplicity, this model captures the
essential physics of confinement and provides a clean derivation of
energy quantization using only the 1D Time-Dependent Schr\"odinger
Equation, the method of separation of variables, and Fourier series
representation.

This report develops that derivation step by step: starting from the
governing partial differential equation, separating it into space and
time parts, applying the physical boundary conditions of the well,
extracting the quantized energy levels, normalizing the resulting
eigenfunctions, and finally assembling the general time-dependent wave
function as a Fourier series expansion.

\section{Objective}
To analytically derive the quantized energy states and formulate the
general wave function $\Psi(x,t)$, expressed as a Fourier series, for an
electron trapped in a one-dimensional infinite potential well. The
derivation is governed by the 1D Time-Dependent Schr\"odinger Equation
and is obtained using the method of separation of variables.

\section{Governing Equation and Setup}
The state of a quantum particle is completely described by its wave
function $\Psi(x,t)$, whose squared magnitude $|\Psi(x,t)|^2$ gives the
probability density of finding the particle at position $x$ at time
$t$. The evolution of this wave function in time is governed by the
1D Time-Dependent Schr\"odinger Equation:
\begin{equation}
    i\hbar \, \frac{\partial \Psi(x,t)}{\partial t}
    = -\frac{\hbar^2}{2m} \, \frac{\partial^2 \Psi(x,t)}{\partial x^2}
    \label{eq:schrodinger}
\end{equation}
Here $\hbar$ is the reduced Planck constant, $m$ is the mass of the
particle (the electron, in this problem), and the right-hand side is
the kinetic energy operator acting on $\Psi$.

For a particle confined to a 1D infinite potential well of length $L$,
the potential energy function is
\begin{equation}
    V(x) =
    \begin{cases}
        0, & 0 \le x \le L \\
        \infty, & \text{otherwise}
    \end{cases}
    \label{eq:potential}
\end{equation}
Inside the well the particle is completely free (no force acts on it),
while outside the well the potential is infinite, meaning it would
take infinite energy for the particle to be found there. Physically,
this means the probability of finding the particle outside the well is
exactly zero, and by continuity of the wave function this forces
$\Psi$ itself to vanish at the walls. This gives the Dirichlet
boundary conditions:
\begin{equation}
    \Psi(0,t) = 0 \qquad \text{and} \qquad \Psi(L,t) = 0 \qquad \forall\, t \ge 0
    \label{eq:bc}
\end{equation}
These two conditions are what will later force the electron's energy
to take only discrete, quantized values.

\section{Method of Separation of Variables}
Equation \eqref{eq:schrodinger} is a linear partial differential
equation in two independent variables, $x$ and $t$. A standard and
powerful technique for solving such linear PDEs -- and the one
required by this assignment -- is separation of variables, where we
look for solutions that factor into a purely spatial part and a purely
temporal part:
\begin{equation}
    \Psi(x,t) = \psi(x)\,\phi(t)
    \label{eq:product}
\end{equation}
This assumption is justified because the boundary conditions in
Equation \eqref{eq:bc} apply only at fixed positions ($x=0,L$) for all
time, which is exactly the situation in which a separable, or
``stationary state,'' solution exists.

Substituting Equation \eqref{eq:product} into Equation
\eqref{eq:schrodinger} and using the product rule (noting that
$\partial \Psi/\partial t = \psi(x)\,\phi'(t)$ and
$\partial^2 \Psi/\partial x^2 = \psi''(x)\,\phi(t)$):
\begin{equation}
    i\hbar\, \psi(x)\,\frac{d\phi(t)}{dt}
    = -\frac{\hbar^2}{2m}\, \phi(t)\, \frac{d^2\psi(x)}{dx^2}
    \label{eq:substituted}
\end{equation}

Dividing both sides of Equation \eqref{eq:substituted} by the product
$\psi(x)\phi(t)$ separates every term so that all $t$-dependence sits
on the left and all $x$-dependence sits on the right:
\begin{equation}
    i\hbar\, \frac{1}{\phi(t)}\frac{d\phi(t)}{dt}
    = -\frac{\hbar^2}{2m}\, \frac{1}{\psi(x)}\frac{d^2\psi(x)}{dx^2}
    \label{eq:before-const}
\end{equation}
The left-hand side of Equation \eqref{eq:before-const} is a function of
$t$ only, while the right-hand side is a function of $x$ only. Since
$x$ and $t$ are independent variables, the only way an expression in
$t$ alone can equal an expression in $x$ alone, for all $x$ and $t$
simultaneously, is if both sides are equal to one and the same
constant. We call this constant $E$, which will turn out to be the
particle's total energy (its eigenvalue):
\begin{equation}
    i\hbar\, \frac{1}{\phi(t)}\frac{d\phi(t)}{dt}
    = -\frac{\hbar^2}{2m}\, \frac{1}{\psi(x)}\frac{d^2\psi(x)}{dx^2}
    = E
    \label{eq:separated}
\end{equation}
This is the elimination-of-variables step: the single PDE in
Equation \eqref{eq:schrodinger} has now been converted into two much
simpler ordinary differential equations (ODEs) -- one purely in $t$
and one purely in $x$ -- linked together only by the shared constant
$E$.

\section{Temporal Solution}
Isolating the time-dependent equality from Equation
\eqref{eq:separated}:
\begin{equation}
    i\hbar\, \frac{d\phi(t)}{dt} = E\,\phi(t)
    \quad \Longrightarrow \quad
    \frac{d\phi(t)}{dt} = -\frac{iE}{\hbar}\,\phi(t)
    \label{eq:temporal-ode}
\end{equation}
Equation \eqref{eq:temporal-ode} is a first-order linear ODE of the
standard form $\phi'(t) = \lambda \phi(t)$, whose solution is a simple
exponential. Solving it gives
\begin{equation}
    \phi(t) = C \exp\!\left(-\frac{iEt}{\hbar}\right)
    \label{eq:temporal}
\end{equation}
where $C$ is an integration constant that will later be absorbed into
the normalization of the spatial part. Because the exponent is purely
imaginary, $\phi(t)$ has constant magnitude $|C|$ and only rotates in
phase as time progresses -- this is why solutions of this form are
called \emph{stationary states}: the probability density
$|\Psi(x,t)|^2 = |\psi(x)|^2 |\phi(t)|^2 = |\psi(x)|^2 |C|^2$ does not
change with time, even though the wave function's phase does.

\section{Spatial Solution and Boundary Conditions}
Isolating the space-dependent equality from Equation
\eqref{eq:separated}:
\begin{equation}
    -\frac{\hbar^2}{2m}\, \frac{d^2\psi(x)}{dx^2} = E\,\psi(x)
    \quad \Longrightarrow \quad
    \frac{d^2\psi(x)}{dx^2} + k^2 \psi(x) = 0,
    \qquad k = \frac{\sqrt{2mE}}{\hbar}
    \label{eq:spatial-ode}
\end{equation}
Equation \eqref{eq:spatial-ode} has exactly the same form as the
classical simple-harmonic-oscillator equation, and so its general
solution is a combination of sine and cosine:
\begin{equation}
    \psi(x) = A\sin(kx) + B\cos(kx)
    \label{eq:spatial-general}
\end{equation}
The constant $k$ here plays the role of a spatial wavenumber, related
to the particle's momentum by $p = \hbar k$.

\subsection*{Applying the Boundary Conditions}
We now use the physical boundary conditions from Equation
\eqref{eq:bc} to pin down the constants $A$, $B$, and the allowed
values of $k$.

\textbf{At $x = 0$:}
\begin{equation}
    \psi(0) = A\sin(0) + B\cos(0) = A(0) + B(1) = B = 0
    \label{eq:bc-left}
\end{equation}
Since $\sin(0)=0$ and $\cos(0)=1$, Equation \eqref{eq:bc-left} forces
$B=0$ directly, leaving the simplified spatial solution
$\psi(x) = A\sin(kx)$.

\textbf{At $x = L$:}
\begin{equation}
    \psi(L) = A\sin(kL) = 0
    \label{eq:bc-right}
\end{equation}
There are two ways to satisfy Equation \eqref{eq:bc-right}: either
$A = 0$, which is the trivial (physically meaningless, ``no particle
at all'') solution, or $\sin(kL) = 0$ with $A \neq 0$. We keep the
non-trivial option, which requires
\begin{equation}
    \sin(kL) = 0 \quad \Longrightarrow \quad k_n L = n\pi,
    \qquad n = 1, 2, 3, \dots
    \label{eq:quantization-condition}
\end{equation}
(the case $n=0$ is excluded because it forces $\psi(x)\equiv 0$
everywhere, again describing no particle). Solving Equation
\eqref{eq:quantization-condition} for $k$ gives the allowed,
discrete wavenumbers:
\begin{equation}
    \boxed{k_n = \dfrac{n\pi}{L}}
    \label{eq:kn}
\end{equation}
This is the heart of quantum confinement: the boundary conditions
permit only standing waves for which a whole number of half-wavelengths
fit exactly inside the well, exactly analogous to the allowed standing
wave modes on a guitar string fixed at both ends.

\section{Derivation of Quantized Energy States}
Having restricted $k$ to the discrete values $k_n$, we now translate
this restriction back into a restriction on energy. Substituting
$k_n = \dfrac{n\pi}{L}$ into the definition $k = \dfrac{\sqrt{2mE}}{\hbar}$
from Equation \eqref{eq:spatial-ode}:
\begin{equation}
    \frac{n\pi}{L} = \frac{\sqrt{2mE_n}}{\hbar}
    \label{eq:k-to-E}
\end{equation}
Squaring both sides of Equation \eqref{eq:k-to-E} removes the square
root:
\begin{equation}
    \frac{n^2\pi^2}{L^2} = \frac{2mE_n}{\hbar^2}
\end{equation}
and solving for $E_n$ gives the quantized energy levels of the
particle in a box:
\begin{resultbox}
\begin{equation}
    E_n = \frac{n^2 \pi^2 \hbar^2}{2mL^2}, \qquad n = 1, 2, 3, \dots
    \label{eq:energy}
\end{equation}
\end{resultbox}
This is a striking departure from classical mechanics: a classical
particle bouncing between two walls could have any energy whatsoever,
but a quantum particle confined the same way can only occupy this
discrete ``ladder'' of energies, spaced further and further apart as
$n$ increases (since $E_n \propto n^2$). Note also that $E_n$ grows as
$1/L^2$ -- the more tightly the electron is confined, the larger the
spacing between its allowed energy levels becomes. This inverse-square
dependence on confinement length is precisely why quantum confinement
effects become significant only at the nanoscale.

\section{Normalization of the Spatial Eigenmodes}
The wave function must satisfy one more physical requirement: since
$|\psi_n(x)|^2$ represents a probability density, the total
probability of finding the particle \emph{somewhere} inside the well
must equal exactly 1. This normalization condition fixes the
remaining constant $A$:
\begin{equation}
    \int_0^L |\psi_n(x)|^2\, dx = 1
    \quad \Longrightarrow \quad
    A^2 \int_0^L \sin^2\!\left(\frac{n\pi x}{L}\right) dx = 1
    \label{eq:norm-integral}
\end{equation}
To evaluate the integral in Equation \eqref{eq:norm-integral}, we use
the half-angle identity $\sin^2\theta = \dfrac{1-\cos 2\theta}{2}$:
\begin{equation}
    \int_0^L \sin^2\!\left(\frac{n\pi x}{L}\right) dx
    = \int_0^L \frac{1}{2}\left[1 - \cos\!\left(\frac{2n\pi x}{L}\right)\right] dx
    = \frac{L}{2} - \underbrace{\frac{L}{4n\pi}\sin\!\left(\frac{2n\pi x}{L}\right)\Big|_0^L}_{=\,0}
    = \frac{L}{2}
\end{equation}
The cosine term integrates to a sine that vanishes at both $x=0$ and
$x=L$ for every integer $n$, leaving simply $L/2$. Substituting this
back into Equation \eqref{eq:norm-integral}:
\begin{equation}
    A^2 \cdot \frac{L}{2} = 1
    \quad \Longrightarrow \quad
    A = \sqrt{\frac{2}{L}}
\end{equation}
Hence the normalized spatial eigenfunctions are
\begin{equation}
    \psi_n(x) = \sqrt{\frac{2}{L}}\, \sin\!\left(\frac{n\pi x}{L}\right)
    \label{eq:psin}
\end{equation}
These functions form an orthonormal set on $[0,L]$, meaning
$\int_0^L \psi_n(x)\psi_m(x)\,dx = 0$ for $n \neq m$ and $=1$ for
$n=m$ -- a property that will be used directly in the next section to
extract the Fourier coefficients.

\section{Final Fourier Series Solution}
Combining the normalized spatial eigenfunction, Equation
\eqref{eq:psin}, with the temporal solution, Equation
\eqref{eq:temporal} (absorbing the constant $C$ into $A$), gives the
complete stationary-state (standing-wave) solution for a single mode
$n$:
\begin{equation}
    \Psi_n(x,t) = \sqrt{\frac{2}{L}}\, \sin\!\left(\frac{n\pi x}{L}\right)
    \exp\!\left(-i\,\frac{E_n t}{\hbar}\right)
    \label{eq:mode-n}
\end{equation}
Each $\Psi_n(x,t)$ individually satisfies the Schr\"odinger equation
and both boundary conditions. However, the electron is not restricted
to occupying only one of these modes; because Equation
\eqref{eq:schrodinger} is linear, \emph{any} linear combination of
solutions is itself a solution. This is the superposition principle,
and it lets us build the most general possible wave function as a sum
over all the allowed modes -- precisely a Fourier series in the
spatial sine-basis $\{\psi_n(x)\}$:
\begin{resultbox}
\begin{equation}
    \Psi(x,t) = \sum_{n=1}^{\infty} c_n \sqrt{\frac{2}{L}}\,
    \sin\!\left(\frac{n\pi x}{L}\right)
    \exp\!\left(-i\, \frac{n^2\pi^2\hbar}{2mL^2}\, t\right)
    \label{eq:general-solution}
\end{equation}
\end{resultbox}

The Fourier coefficients $c_n$ are not free parameters; they are
completely determined by the particle's initial wave function
$\Psi(x,0)$ at $t=0$. Setting $t=0$ in Equation
\eqref{eq:general-solution} gives
$\Psi(x,0) = \sum_n c_n \psi_n(x)$, and because the $\psi_n(x)$ are
orthonormal, we can project out any single coefficient by multiplying
both sides by $\psi_n(x)$ and integrating over $[0,L]$:
\begin{equation}
    c_n = \int_0^L \Psi(x,0)\, \sqrt{\frac{2}{L}}\, \sin\!\left(\frac{n\pi x}{L}\right) dx
    \label{eq:cn}
\end{equation}

\subsection*{Physical Interpretation}
Equation \eqref{eq:general-solution} states that \emph{any} allowed
wave function inside the box, however complicated, can be written as a
weighted sum of the discrete standing-wave modes $\psi_n(x)$, each one
oscillating in time at its own characteristic frequency $E_n/\hbar$.
The weights $c_n$ are set entirely by the initial condition through
Equation \eqref{eq:cn}. This is exactly the same eigenfunction-expansion
idea used to build an ordinary Fourier series on an interval -- only
here the ``basis'' functions are not a generic sine/cosine set but the
specific quantized eigenmodes forced upon the system by the physical
boundaries of the potential well.

\section{Conclusion}
Starting from the 1D Time-Dependent Schr\"odinger Equation, this
report has shown, step by step, how confining an electron to a region
of length $L$ with infinitely high walls leads directly to
quantization. Separation of variables split the governing PDE into an
exponential temporal equation and a harmonic-oscillator-type spatial
equation. Applying the physical boundary conditions
$\Psi(0,t)=\Psi(L,t)=0$ restricted the spatial solutions to a discrete
family of sine modes, $\psi_n(x) = \sqrt{2/L}\,\sin(n\pi x/L)$, and
correspondingly restricted the particle's energy to the discrete
ladder $E_n = n^2\pi^2\hbar^2/(2mL^2)$.

Combining the normalized spatial modes with their time evolution and
invoking the superposition principle produced the general solution as
a Fourier series, Equation \eqref{eq:general-solution}, with
coefficients $c_n$ fixed by the initial state via Equation
\eqref{eq:cn}. Two engineering conclusions follow directly from this
derivation. First, confinement is what \emph{creates} discreteness --
a free electron ($L \to \infty$) has a continuous energy spectrum, but
as $L$ shrinks the energy levels spread apart as $1/L^2$, which is
precisely the mechanism exploited in quantum dots, quantum wells, and
nanoscale transistor channels to engineer specific, tunable electronic
and optical properties. Second, the Fourier series structure of the
general solution shows that the time evolution of \emph{any} initial
electron state in such a structure can be predicted exactly, simply by
decomposing it into these quantized eigenmodes -- a technique that
extends naturally to more realistic 2D and 3D confinement geometries
used in real semiconductor device design.

\end{document}
